\documentclass[aps,prc,onecolumn,superscriptaddress,nofootinbib,10pt]{revtex4-2}
\usepackage{amsmath,amssymb}
\usepackage{booktabs}
\usepackage[colorlinks=true,linkcolor=blue,citecolor=blue,urlcolor=blue]{hyperref}

\newcommand{\nB}{n_B}
\newcommand{\nsat}{n_{\mathrm{sat}}}
\newcommand{\muB}{\mu_B}
\newcommand{\muC}{\mu_C}
\newcommand{\muS}{\mu_S}
\newcommand{\munue}{\mu_{\nu_e}}
\newcommand{\SigR}{\Sigma^{R}}
\newcommand{\mueff}[1]{\mu^{\mathrm{eff}}_{#1}}
\newcommand{\meff}[1]{m^{*}_{#1}}
\newcommand{\hc}{(\hbar c)^3}

\begin{document}

\title{DD2: a density-dependent relativistic mean-field model \\
       for hadronic matter --- as implemented in \texttt{eos/dd2}}

\author{M.~Guerrini}
\affiliation{University of Ferrara}

\begin{abstract}
The \texttt{eos/dd2} subpackage implements the DD2 parametrization of the
density-dependent relativistic mean-field (DD-RMF) model of Typel et
al.~\cite{Typel2010}, extended to the hyperon octet with the DD2Y couplings of
Marques et al.~\cite{Marques2017,Fortin2017}, to the $\Delta(1232)$ quartet,
and to a thermal pseudoscalar/vector meson gas following
Lavagno~\cite{Lavagno2010} (see also arXiv:1210.0400). This note states the
Lagrangian, the parameter set, the field equations, the thermodynamics
including the rearrangement self-energy, and the closure of each equilibrium
mode the code supports. It is written from the code, as its specification:
every equation here is asserted by the \texttt{verify/} suite or the test
gates.
\end{abstract}

\maketitle

\section{Lagrangian and couplings}

Baryons $i$ (nucleons; optionally the hyperon octet and the $\Delta$ quartet)
couple to the isoscalar--scalar $\sigma$, isoscalar--vector $\omega$,
isovector--vector $\rho$, and (for strange baryons) the hidden-strange vector
$\phi$:
%
\begin{align}
\mathcal{L} ={}& \sum_i \bar\psi_i \Big[ \gamma_\mu \big( i\partial^\mu
      - \Gamma_{\omega i}\,\omega^\mu
      - \Gamma_{\rho i}\, t_{3i}\, \rho^\mu
      - \Gamma_{\phi i}\, \phi^\mu \big)
      - \big( m_i - \Gamma_{\sigma i}\,\sigma \big) \Big] \psi_i
      \nonumber\\
 &+ \tfrac12 \big( \partial_\mu \sigma\, \partial^\mu \sigma
      - m_\sigma^2 \sigma^2 \big)
   - \tfrac14 W_{\mu\nu} W^{\mu\nu} + \tfrac12 m_\omega^2\, \omega_\mu \omega^\mu
   - \tfrac14 R_{\mu\nu} R^{\mu\nu} + \tfrac12 m_\rho^2\, \rho_\mu \rho^\mu
   - \tfrac14 \Phi_{\mu\nu} \Phi^{\mu\nu} + \tfrac12 m_\phi^2\, \phi_\mu \phi^\mu ,
\label{eq:lagrangian}
\end{align}
%
with $t_{3i}$ the third isospin component in the $\tau_3 = \pm 1$ convention
($t_{3p} = +1$, $t_{3n} = -1$), so the published DD2 $\rho$ coupling is used
as tabulated. Leptons ($e$, $\mu$) and, in trapped modes, electron neutrinos
enter as free Dirac gases; photons as a free Bose gas.

The nucleon couplings depend on the total baryon density,
$\Gamma_{iN}(\nB) = \Gamma_{iN}(\nsat) f_i(x)$ with $x = \nB/\nsat$, in the
Typel--Wolter form~\cite{TypelWolter1999,Typel2010}:
%
\begin{align}
f_i(x) &= a_i \frac{1 + b_i (x + d_i)^2}{1 + c_i (x + d_i)^2},
  \qquad i = \sigma, \omega,
\label{eq:frational}\\
f_\rho(x) &= \exp\!\big[ -a_\rho (x - 1) \big].
\label{eq:fexp}
\end{align}
%
Two structural constraints make two of the four isoscalar shape coefficients
dependent,
%
\begin{equation}
f_i(1) = 1 \;\Rightarrow\; a_i = \frac{1 + c_i (1+d_i)^2}{1 + b_i (1+d_i)^2},
\qquad
f_i''(0) = 0 \;\Rightarrow\; d_i = \frac{1}{\sqrt{3 c_i}} ,
\label{eq:constraints}
\end{equation}
%
and every ingested parameter set is validated against
Eq.~\eqref{eq:constraints} to $2\times 10^{-6}$ (the accuracy to which the
published table itself obeys them). The published DD2 values
\cite{Typel2010} are transcribed verbatim in
Table~\ref{tab:params}.

\begin{table}[b]
\caption{The DD2 parameter set (Typel et al.~\cite{Typel2010}), as carried by
\texttt{Parameters.default()}. $\nsat$ doubles as the
coupling reference density. The $\phi$ mass ($1019.45$~MeV) enters only with
hyperons.}
\label{tab:params}
\begin{ruledtabular}
\begin{tabular}{lcc|lcc}
$\nsat$ & $0.149065$ & fm$^{-3}$ & $\Gamma_\sigma(\nsat)$ & $10.686681$ & \\
$m_n$   & $939.56536$ & MeV      & $a_\sigma$ & $1.357630$ & \\
$m_p$   & $938.27203$ & MeV      & $b_\sigma$ & $0.634442$ & \\
$m_\sigma$ & $546.212459$ & MeV  & $c_\sigma$ & $1.005358$ & \\
$m_\omega$ & $783.0$ & MeV       & $d_\sigma$ & $0.575810$ & \\
$m_\rho$   & $763.0$ & MeV       & $\Gamma_\omega(\nsat)$ & $13.342362$ & \\
 & & & $a_\omega$ & $1.369718$ & \\
 & & & $b_\omega$ & $0.496475$ & \\
 & & & $c_\omega$ & $0.817753$ & \\
 & & & $d_\omega$ & $0.638452$ & \\
 & & & $\Gamma_\rho(\nsat)$ & $3.626940$ & \\
 & & & $a_\rho$ & $0.518903$ & \\
\end{tabular}
\end{ruledtabular}
\end{table}

Non-nucleonic baryons inherit the nucleon density dependence through constant
ratios $x_{iY} = \Gamma_{iY}/\Gamma_{iN}$. The vector ratios are SU(6):
$x_{\omega\Lambda} = x_{\omega\Sigma} = 2/3$, $x_{\omega\Xi} = 1/3$;
$x_{\rho\Lambda} = 0$, $x_{\rho\Sigma} = 2$, $x_{\rho\Xi} = 1$; and the
hidden-strange vector couples through
$g_{\phi Y}/g_{\omega N} = \{-\sqrt2/3,\, -\sqrt2/3,\, -2\sqrt2/3\}$ for
$\{\Lambda, \Sigma, \Xi\}$, inheriting the $\omega$ density dependence:
$\Gamma_{\phi Y}(\nB) = x_{\phi Y}\, \Gamma_{\omega N}(\nB)$ (the DD2Y
choice~\cite{Marques2017}). The scalar ratios are either the published DD2Y
values $R_\sigma = \{0.62, 0.48, 0.32\}$~\cite{Fortin2017} or inverted from
the hyperon single-particle potentials in symmetric matter at saturation,
%
\begin{equation}
U_Y(\nsat) = -\Gamma_{\sigma Y}\, \bar\sigma
             + \Gamma_{\omega Y}\, \omega_0 + \SigR
\qquad (\text{defaults } U_\Lambda = -30,\; U_\Sigma = +30,\;
        U_\Xi = -18 \text{ MeV}),
\label{eq:UY}
\end{equation}
%
one linear equation per hyperon; with $U_\Xi = -18$~MeV this inversion
regenerates the published $R_\sigma$ column. The $\Delta$ quartet has no
canonical DD2 coupling table: the default is universal coupling
($x_{\Delta i} = 1$), or $x_{\Delta\sigma}$ is fixed by the same inversion
from a chosen $U_\Delta \in [-100, -50]$~MeV.

\section{Mean fields, rearrangement, and thermodynamics}

In static uniform matter the meson field equations are algebraic,
%
\begin{equation}
m_\sigma^2 \sigma = \sum_i \Gamma_{\sigma i}\, n^s_i, \qquad
m_\omega^2 \omega_0 = \sum_i \Gamma_{\omega i}\, n_i, \qquad
m_\rho^2 \rho_0 = \sum_i \Gamma_{\rho i}\, t_{3i}\, n_i, \qquad
m_\phi^2 \phi_0 = \sum_i \Gamma_{\phi i}\, n_i ,
\label{eq:fieldeqs}
\end{equation}
%
with $n_i$ and $n^s_i$ the number and scalar densities of an ideal Fermi gas
of effective mass and effective chemical potential
%
\begin{equation}
\meff{i} = m_i - \Gamma_{\sigma i}\, \sigma, \qquad
\mueff{i} = \mu_i - \Gamma_{\omega i}\, \omega_0
          - \Gamma_{\rho i}\, t_{3i}\, \rho_0
          - \Gamma_{\phi i}\, \phi_0 - \SigR .
\label{eq:mueff}
\end{equation}
%
In code, $\mueff{i}$ and $\meff{i}$ are named \texttt{mu\_eff\_i} and
\texttt{m\_eff\_i}.

\subsection{One species as an ideal gas}
\label{sec:kinetic}

Every baryon enters as an ideal Fermi gas of mass $\meff{i}$ at potential
$\mueff{i}$, with degeneracy $g_i$ ($2$ for the octet, $4$ for the $\Delta$)
and antiparticles included. With $E = \sqrt{k^2 + \meff{i}{}^2}$ and
$f_\pm = \big[ 1 + e^{(E \mp \mueff{i})/T} \big]^{-1}$,
%
\begin{align}
n_i &= \frac{g_i}{2\pi^2 (\hbar c)^3} \int_0^\infty \!\! dk\; k^2
       \big( f_+ - f_- \big) ,
\qquad
\varepsilon^{\mathrm{kin}}_i = \frac{g_i}{2\pi^2 (\hbar c)^3}
       \int_0^\infty \!\! dk\; k^2 E \big( f_+ + f_- \big) ,
\nonumber\\
P^{\mathrm{kin}}_i &= \frac{g_i}{6\pi^2 (\hbar c)^3} \int_0^\infty \!\! dk\;
       \frac{k^4}{E} \big( f_+ + f_- \big) .
\label{eq:kin}
\end{align}
%
The scalar density and the entropy density are \emph{not} integrated. They
follow from the trace of the energy--momentum tensor and from the Euler
relation for a single species,
%
\begin{equation}
n^s_i = \frac{\varepsilon^{\mathrm{kin}}_i - 3 P^{\mathrm{kin}}_i}{\meff{i}} ,
\qquad
s_i = \frac{\varepsilon^{\mathrm{kin}}_i + P^{\mathrm{kin}}_i
            - \mueff{i}\, n_i}{T} ,
\label{eq:nsandS}
\end{equation}
%
which is how the code obtains them, and is worth stating because it means an
error in $\varepsilon^{\mathrm{kin}}_i$ or $P^{\mathrm{kin}}_i$ propagates
into $n^s_i$ --- and $n^s_i$ sources the $\sigma$ field, so such an error does
not stay confined to the totals.

At $T = 0$ the integrals are elementary. With
$k_F = \sqrt{\mueff{i}{}^2 - \meff{i}{}^2}$ (everything vanishing when
$|\mueff{i}| \le \meff{i}$) and $L = \ln\big[(k_F + |\mueff{i}|)/\meff{i}\big]$,
%
\begin{align}
n_i &= \mathrm{sgn}(\mueff{i})\, \frac{g_i\, k_F^3}{6\pi^2 (\hbar c)^3} ,
\qquad
n^s_i = \frac{g_i\, \meff{i}}{4\pi^2 (\hbar c)^3}
        \Big( k_F |\mueff{i}| - \meff{i}{}^2 L \Big) ,
\nonumber\\
P^{\mathrm{kin}}_i &= \frac{g_i}{48\pi^2 (\hbar c)^3}
       \Big[ \big( 2k_F^3 - 3\meff{i}{}^2 k_F \big) |\mueff{i}|
             + 3 \meff{i}{}^4 L \Big] ,
\label{eq:T0}\\
\varepsilon^{\mathrm{kin}}_i &= \frac{g_i}{16\pi^2 (\hbar c)^3}
       \Big[ \big( 2k_F^3 + \meff{i}{}^2 k_F \big) |\mueff{i}|
             - \meff{i}{}^4 L \Big] ,
\qquad s_i = 0 .
\nonumber
\end{align}
%
At $T>0$ Eq.~\eqref{eq:kin} is evaluated with the Johns--Ellis--Lattimer
rational approximants~\cite{JohnsEllisLattimer1996}, accurate to
$\sim 10^{-4}$, from \texttt{eos/general/fermi\_integrals} --- the single home
of all Fermi and Bose integrals in this repository, so that no model carries
its own. A Gauss--Laguerre quadrature of the same integrals is available there
as the accuracy reference. The $T = 0$ branch is additionally Numba-compiled
(\S\,\ref{sec:numerics}).

The conserved-charge densities assembled from these are
%
\begin{equation}
\nB = \sum_i B_i n_i , \qquad
n_C = \sum_i Q_i n_i + n_C^{\mathrm{mes}} , \qquad
n_S = \sum_i S_i n_i + n_S^{\mathrm{mes}} ,
\label{eq:charges}
\end{equation}
%
the meson terms being those of Sec.~\ref{sec:gas}. Note that $n_C$ and $n_S$
are the TOTALS, gas included: that is what the neutrality and fixed-fraction
conditions below are stated in terms of.

The density dependence of the couplings produces the rearrangement
self-energy, identical for all baryons because the couplings depend on the
total $\nB$:
%
\begin{equation}
\SigR = \frac{\partial \Gamma_{\omega N}}{\partial \nB}\, \omega_0
        \sum_i x_{\omega i} n_i
      + \frac{\partial \Gamma_{\rho N}}{\partial \nB}\, \rho_0
        \sum_i x_{\rho i} t_{3i} n_i
      + \frac{\partial \Gamma_{\omega N}}{\partial \nB}\, \phi_0
        \sum_i x_{\phi i} n_i
      - \frac{\partial \Gamma_{\sigma N}}{\partial \nB}\, \sigma
        \sum_i x_{\sigma i} n^s_i .
\label{eq:sigmaR}
\end{equation}
%
$\SigR$ enters the chemical potentials and the pressure, never the energy
density --- this is what makes the model thermodynamically consistent:
%
\begin{align}
\varepsilon &= \sum_i \varepsilon^{\mathrm{kin}}_i
  + \tfrac12 \big( m_\sigma^2 \sigma^2 + m_\omega^2 \omega_0^2
                 + m_\rho^2 \rho_0^2 + m_\phi^2 \phi_0^2 \big)
  + \varepsilon_{\mathrm{lep}} + \varepsilon_\gamma + \varepsilon_{\mathrm{mes}},
\label{eq:eps}\\
P &= \sum_i P^{\mathrm{kin}}_i
  + \tfrac12 \big( -m_\sigma^2 \sigma^2 + m_\omega^2 \omega_0^2
                 + m_\rho^2 \rho_0^2 + m_\phi^2 \phi_0^2 \big)
  + \nB\, \SigR + P_{\mathrm{lep}} + P_\gamma + P_{\mathrm{mes}} ,
\label{eq:P}\\
s &= \sum_i s_i + s_{\mathrm{lep}} + s_\gamma + s_{\mathrm{mes}} .
\label{eq:s}
\end{align}
%
The mean fields carry no entropy, so $s$ has no field term --- the asymmetry
between Eqs.~\eqref{eq:eps} and \eqref{eq:P} is entirely the sign of the
$\sigma$ mass term and the rearrangement term $\nB \SigR$, and those are the
two places a mean-field model is most easily got wrong. Leptons and photons
are added by the solver rather than by the matter block, because they are
shared by the whole system rather than belonging to a phase: electrons and
muons through Eqs.~\eqref{eq:kin}--\eqref{eq:T0} at their own masses with
$g = 2$, trapped electron neutrinos likewise at $m = 0$, $g = 1$, and photons
as
%
\begin{equation}
P_\gamma = \frac{\pi^2 T^4}{45 (\hbar c)^3}, \qquad
\varepsilon_\gamma = 3 P_\gamma, \qquad
s_\gamma = \frac{4\pi^2 T^3}{45 (\hbar c)^3} .
\label{eq:photons}
\end{equation}
%
Every assembled state is checked against the Hugenholtz--Van Hove
identity~\cite{HugenholtzVanHove1958}
$\varepsilon + P - T s = \sum_i \mu_i n_i$ to a relative $10^{-8}$; a
violation raises rather than returning a state. The sum runs over every
species present, with the baryons and leptons at their FULL potentials and the
thermal meson gas at its EFFECTIVE ones,
%
\begin{equation}
\sum_i \mu_i n_i
  = \sum_i \big( B_i \muB + Q_i \muC + S_i \muS \big) n_i
  + \sum_j \mu^{*}_j n_j
  + \mu_e n_e + \mu_\mu n_\mu + \munue n_{\nu_e} ,
\label{eq:eulersum}
\end{equation}
%
because the vector shifts the gas carries are sourced by the baryons alone and
so have no partner in the field energy to cancel against.

The nucleon-mass convention of the uniform-matter kernel is selectable: the
default (\texttt{"average"}) gives both nucleons the mass
$(m_n + m_p)/2$, which is the convention the published table and the golden
regression values are stated in; \texttt{"physical"} uses per-species masses,
needed for sub-$0.1\%$ isovector agreement with the HS(DD2) CompOSE
tables~\cite{HempelSchaffnerBielich2010,TypelCompOSE2015}.

\section{Conserved charges and species potentials}

The conserved-charge basis is $(B, C, S)$ with $C$ the electric charge of
strongly interacting matter only and $S = +1$ per strange quark ($\Lambda$:
$S=+1$, $\Xi$: $S=+2$; the opposite of the PDG sign, used consistently).
Species potentials are derived, never independent unknowns:
%
\begin{equation}
\mu_i = B_i\, \muB + Q_i\, \muC + S_i\, \muS ,
\qquad\text{so}\qquad
\muC = \mu_p - \mu_n,
\quad \text{and beta equilibrium reads } \muC + \mu_e = 0 .
\label{eq:basis}
\end{equation}
%
In trapped-neutrino matter the electron family adds $\mu_{L}$:
$\mu_e = \munue - \muC$, $\mu_{\nu_e} = \munue$; the muon family remains
transparent ($\mu_\mu = -\muC$, no $\nu_\mu$ tracked).

\section{The thermal meson gas}
\label{sec:gas}

At $T>0$ the pseudoscalar nonet ($\pi$, $K$, $\eta$, $\eta'$) and optionally
the vector nonet ($\rho$, $\omega$, $K^*$, $\phi$) ride on the mean field as
one-body Bose gases~\cite{Lavagno2010} (see also arXiv:1210.0400), with
effective potentials tied to the same fields that generate the baryon
interactions, using DD2's density-dependent couplings in place of constant
ones:
%
\begin{align}
\mu^{*}_{\pi^+} = \mu^{*}_{\rho^+} &= \muC - \Gamma_{\rho N}(\nB)\, \rho_0 ,
\nonumber\\
\mu^{*}_{K^+} = \mu^{*}_{K^{*+}} &= \muC - \muS
   - \big[ \Gamma_{\omega N} - \Gamma_{\omega \Lambda} \big](\nB)\, \omega_0
   - \tfrac12\, \Gamma_{\rho N}(\nB)\, \rho_0 ,
\label{eq:mesonshifts}\\
\mu^{*}_{K^0} = \mu^{*}_{K^{*0}} &= - \muS
   - \big[ \Gamma_{\omega N} - \Gamma_{\omega \Lambda} \big](\nB)\, \omega_0
   + \tfrac12\, \Gamma_{\rho N}(\nB)\, \rho_0 ,
\nonumber
\end{align}
%
with $\mu^{*}_{j^-} = -\mu^{*}_{j^+}$ and $\mu^{*}_j = 0$ for strangeless
neutral mesons. No rearrangement term enters $\mu^{*}_j$: the gas is a
spectator to $\SigR$. The gas carries electric charge and strangeness ---
under this sign convention $K^-$ and $\bar K^0$ carry $S=+1$ --- and those
densities \emph{enter the equilibrium constraints}: charge neutrality and the
fixed-$Y_C$/$Y_S$ conditions are imposed on baryons and mesons together. The
gas carries no baryon number and does not source the field equations.

\subsection{The Bose thermodynamics of one meson species}

Each species is an ideal Bose gas of mass $m_j$, degeneracy $g_j$ and
effective potential $\mu^*_j$, with its antiparticle carried as a
\emph{separate} species at $-\mu^*_j$ rather than through an antiparticle term
inside the integral:
%
\begin{align}
n_j &= \frac{g_j}{2\pi^2 \hc} \int_0^\infty \! dk\, k^2\, b(E_k - \mu^*_j) ,
      & E_k &= \sqrt{k^2 + m_j^2} ,
\label{eq:nmeson} \\
P_j &= \frac{g_j}{6\pi^2 \hc} \int_0^\infty \! dk\, \frac{k^4}{E_k}\,
       b(E_k - \mu^*_j) ,
\label{eq:Pmeson} \\
\varepsilon_j &= \frac{g_j}{2\pi^2 \hc} \int_0^\infty \! dk\, k^2 E_k\,
       b(E_k - \mu^*_j) ,
\label{eq:epsmeson} \\
s_j &= \frac{\varepsilon_j + P_j - \mu^*_j n_j}{T} ,
\label{eq:smeson}
\end{align}
%
with the Bose occupation $b(x) = [e^{x/T} - 1]^{-1}$. Eq.~\eqref{eq:smeson} is
the one-species Euler relation again, at the \emph{effective} potential --- the
same identity Eq.~\eqref{eq:nsandS} uses for a baryon, which is why the gas
joins the Hugenholtz--Van Hove sum through $\sum_j \mu^*_j n_j$ and not
through a physical potential. These integrals come from
\texttt{eos.general.bose\_integrals}, evaluated by the same
Johns--Ellis--Lattimer construction as the Fermi ones. There is no $T = 0$
branch: the gas is defined only at $T > 0$ and returns zero below it.

The active species, their masses (MeV), charges, strangeness under this
repository's $S = +1$ per $s$ quark, and degeneracies:
%
\begin{center}
\begin{tabular}{lrrrr@{\qquad}lrrrr}
\toprule
species & $m_j$ & $Q_j$ & $S_j$ & $g_j$ &
species & $m_j$ & $Q_j$ & $S_j$ & $g_j$ \\
\midrule
$\pi^+$    & $139.57039$ & $+1$ & $0$  & $1$ &
$\rho^+$   & $775.26$    & $+1$ & $0$  & $3$ \\
$\pi^-$    & $139.57039$ & $-1$ & $0$  & $1$ &
$\rho^-$   & $775.26$    & $-1$ & $0$  & $3$ \\
$\pi^0$    & $134.9768$  & $0$  & $0$  & $1$ &
$\rho^0$   & $775.26$    & $0$  & $0$  & $3$ \\
$K^+$      & $493.677$   & $+1$ & $-1$ & $1$ &
$\omega$   & $782.66$    & $0$  & $0$  & $3$ \\
$K^-$      & $493.677$   & $-1$ & $+1$ & $1$ &
$K^{*+}$   & $891.67$    & $+1$ & $-1$ & $3$ \\
$K^0$      & $497.611$   & $0$  & $-1$ & $1$ &
$K^{*-}$   & $891.67$    & $-1$ & $+1$ & $3$ \\
$\bar K^0$ & $497.611$   & $0$  & $+1$ & $1$ &
$K^{*0}$   & $891.67$    & $0$  & $-1$ & $3$ \\
$\eta$     & $547.862$   & $0$  & $0$  & $1$ &
$\bar K^{*0}$ & $891.67$ & $0$  & $+1$ & $3$ \\
$\eta'$    & $957.78$    & $0$  & $0$  & $1$ &
$\phi$     & $1019.461$  & $0$  & $0$  & $3$ \\
\bottomrule
\end{tabular}
\end{center}
%
The left column is \texttt{thermal\_mesons}; the right is the optional vector
nonet, which reuses the same three potentials, because the shift depends on
the quark content and vector and pseudoscalar partners share it. Note
$g = 1$ for a pseudoscalar and $g = 3$ for a vector, and that antiparticles
are separate rows: a gas of $\pi^+$ and $\pi^-$ at $\pm\mu^*$ is two species,
not one with an antiparticle term.

\subsection{Condensation is refused, not approximated}

Eqs.~\eqref{eq:nmeson}--\eqref{eq:smeson} describe the gas only while
$|\mu^*_j| < m_j$ for every species. At $|\mu^*_j| = m_j$ the species
condenses: the particles beyond the critical density go into the $k = 0$
state, carrying charge and $\varepsilon = m_j n_{\rm cond}$ with \emph{no}
pressure and \emph{no} entropy, and $n_{\rm cond}$ becomes a new unknown
rather than a function of $\mu^*_j$. None of that is implemented. So the
package computes the diagnostic
%
\begin{equation}
\text{condensation} \;=\; \max_j \frac{|\mu^*_j|}{m_j}
\label{eq:condensation}
\end{equation}
%
and \textbf{refuses any state with condensation $\ge 1$}, raising rather than
returning it. The refusal is not defensive tidiness: the underlying integral
routine \emph{caps} $\mu^*$ at $m$, so past the threshold the gas silently
stops absorbing charge and the entropy it reports drifts --- turning negative
by $|\mu^*|/m \simeq 3$. A returned state there would be wrong rather than
approximate. \texttt{eos.sfho} refuses the same condition, there as a status
because its solvers report rather than raise.

Its $P$, $\varepsilon$, $s$ and $\sum_j \mu^{*}_j n_j$ join the totals and the
Hugenholtz--Van Hove sum.

\section{Equilibrium modes and their closures}

All modes are solved by one residual system over the unknown vector
%
\begin{equation}
x = \big[ \sigma,\ \omega_0,\ \rho_0,\ (\phi_0),\ \tilde\mu_B,\ \muC,\
          (\muS),\ (\munue) \big],
\qquad \tilde\mu_B \equiv \muB - \SigR ,
\label{eq:unknowns}
\end{equation}
%
where $\phi_0$ is present iff hyperons and the $\phi$ field are active,
$\muS$ iff strangeness is fixed, and $\munue$ iff neutrinos are trapped.
Working with $\tilde\mu_B$ (and species $\mueff{i}$) instead of $\muB$
removes the rearrangement circularity and the large vector shift from the
iteration --- $\muB = \tilde\mu_B + \SigR$ is restored at assembly --- and
makes warm-started density sweeps well conditioned, because the effective
potentials vary smoothly along $\nB$ where the full potentials jump.

The residual is assembled in this order, and these are all of its rows. The
mean fields are eliminated against their own sources, so a field row is
$\text{field} - \text{source}/m_M^2$ rather than the raw field equation, and
every row is divided by $m_N$ or by $\nB$ so that all of them are
dimensionless and of order unity:
%
\begin{align}
R_1 &= \Big[ \sigma - \textstyle\sum_i \Gamma_{\sigma i} n^s_i \big/ m_\sigma^2
        \Big] \big/ m_N ,
\qquad
R_2 = \Big[ \omega_0 - \textstyle\sum_i \Gamma_{\omega i} n_i \big/ m_\omega^2
        \Big] \big/ m_N ,
\nonumber\\
R_3 &= \Big[ \rho_0 - \textstyle\sum_i \Gamma_{\rho i} t_{3i} n_i
        \big/ m_\rho^2 \Big] \big/ m_N ,
\qquad
R_4 = \Big[ \phi_0 - \textstyle\sum_i \Gamma_{\phi i} n_i \big/ m_\phi^2
        \Big] \big/ m_N
\;\;\text{(iff $\phi_0$ active)} ,
\nonumber\\
R_5 &= \Big[ \textstyle\sum_i B_i n_i - \nB \Big] \big/ \nB ,
\label{eq:residual}\\
R_6 &= \Big[ n_C - n_e - n_\mu \Big] \big/ \nB
        \quad\text{($C$ equilibrated)} ,
\qquad
R_6 = \Big[ n_C - Y_C \nB \Big] \big/ \nB
        \quad\text{($Y_C$ imposed)} ,
\nonumber\\
R_7 &= \Big[ n_S - Y_S \nB \Big] \big/ \nB
        \qquad \text{(iff $Y_S$ imposed)} ,
\nonumber\\
R_8 &= \Big[ n_e + n_{\nu_e} - Y_{L_e} \nB \Big] \big/ \nB
        \qquad \text{(iff the electron family is trapped)} .
\nonumber
\end{align}
%
$n_C$ and $n_S$ are Eq.~\eqref{eq:charges}'s totals, thermal meson gas
included; the gas carries no baryon number, so it is absent from $R_5$. Per
mode:
%
\begin{center}
\begin{tabular}{llll}
\toprule
mode & independent variables & closing constraints & potentials \\
\midrule
\texttt{beta\_eq\_neutrinoless} & $(\nB, T)$ &
  $\sum_i Q_i n_i + n_C^{\mathrm{mes}} - n_e - n_\mu = 0$ &
  $\muS = \munue = 0$ \\
\texttt{beta\_eq\_neutrino\_trapped} & $(\nB, Y_{L_e}, T)$ &
  neutrality; $(n_e + n_{\nu_e})/\nB = Y_{L_e}$ &
  $\muS = 0$; $\munue$ unknown \\
\texttt{fixed\_YC} & $(\nB, Y_C, T)$ &
  $\big(\textstyle\sum_i Q_i n_i + n_C^{\mathrm{mes}}\big)/\nB = Y_C$ &
  $\muS = \munue = 0$ \\
\texttt{fixed\_YC\_YS} & $(\nB, Y_C, Y_S, T)$ &
  the $Y_C$ row; $\big(\textstyle\sum_i S_i n_i + n_S^{\mathrm{mes}}\big)/\nB = Y_S$ &
  $\muS$ unknown; $\munue = 0$ \\
\bottomrule
\end{tabular}
\end{center}
%
In \texttt{fixed\_YC} the flag \texttt{leptons} selects between the
strongly-interacting slice alone (electrically charged; what a mixed-phase
construction consumes) and the same slice plus neutralizing leptons
($\mu_\mu = \mu_e$, $n_e + n_\mu = Y_C \nB$) --- the leptons do not source
the fields, so they close post hoc through a single 1-D root. Wherever a
temperature is accepted, entropy per baryon $s/\nB$ is accepted in its place
through an outer 1-D solve for $T$ (monotone at fixed $\nB$, so the bracket
is well posed).

\subsection{The reduced nucleon-only system}
\label{sec:reduced}

For nucleons-only beta equilibrium the package carries a \emph{second} system,
with its own unknown vector, because eliminating $\omega_0$ collapses two
unknowns into none: $x = [\sigma, \rho_0, \mueff{n}, \muC]$, four unknowns and
four rows. $\omega_0$ is eliminated algebraically at the target density ---
which works only while every species shares the nucleon couplings, and is
exactly why the octet system above cannot do it --- and
%
\begin{equation}
\mueff{p} = \mueff{n} + \muC
            - \big(\tau_3^{(p)} - \tau_3^{(n)}\big)\, \Gamma_\rho \rho_0
          = \mueff{n} + \muC - 2\,\Gamma_\rho \rho_0
\end{equation}
%
under the $\tau_3 = \pm 1$ convention. The four rows, in the order assembled
and each already dimensionless:
%
\begin{align}
\bar R_1 &= \Big[ \sigma
   - \Gamma_\sigma \big( n^s_n + n^s_p \big) \big/ m_\sigma^2 \Big]
   \big/ \bar m , \\
\bar R_2 &= \Big[ \rho_0
   - \Gamma_\rho \big( \tau_3^{(n)} n_n + \tau_3^{(p)} n_p \big)
     \big/ m_\rho^2 \Big] \big/ \bar m , \\
\bar R_3 &= \frac{n_n + n_p}{\nB} - 1 , \\
\bar R_4 &= \frac{n_p - n_e - n_\mu}{\nB} ,
\end{align}
%
with $\bar m$ the mean nucleon mass, $\mu_e = -\muC$ and
$\mu_\mu = \mu_e$. $\bar R_1$ and $\bar R_2$ are the $\sigma$ and $\rho$ field
equations, $\bar R_3$ baryon number, $\bar R_4$ electric neutrality. The
system is guarded rather than allowed to wander: if either effective mass goes
non-positive the residual returns $[10^6, 0, 0, 0]$, which pushes the root
finder back inside the physical domain instead of letting it evaluate a Fermi
integral at an imaginary mass.

\subsection{The phase-adapter residual}
\label{sec:adapter}

A third system exists, and it is not a mode: it is what \texttt{eos/mixed}
consumes across the phase-adapter contract. Given
$(\tilde\mu_B, \muC, \muS, T)$ --- the potentials as \emph{inputs} --- it
solves for
%
\begin{equation}
x = [\sigma, \omega_0, \rho_0, (\phi_0), \nB] ,
\end{equation}
%
the field equations plus one baryon-density self-consistency row. \textbf{The
density is an unknown here}, unlike in every mode above, because for one phase
of a mixture only the \emph{average} density is prescribed; and because DD2's
couplings are density-dependent, $\Gamma_i(\nB)$ is re-evaluated at the
current $\nB$ on every iteration, which is what makes this a coupled system
rather than a field solve at fixed couplings.

\textbf{There is no charge, strangeness or neutrality row.} That absence is the
content of the contract: the potentials are given, so nothing here chooses a
composition, and \emph{that} is what makes this thermodynamics rather than a
mode. Imposing neutrality on one phase of a Gibbs construction would be wrong
in any case --- global neutrality is the engine's condition, not the phase's.
A non-positive $\nB$ iterate returns a large residual on every row, the same
guard as Sec.~\ref{sec:reduced}.

\section{Nuclear-matter parameters}

The forward map evaluates, at the parametrization's own saturation density
(the $P=0$ root of symmetric matter): $E_{\mathrm{sat}} = E/A$,
$m^*/m$, $K_{\mathrm{sat}} = 9 \nsat^2\, \partial^2_n (E/A)$,
$Q_{\mathrm{sat}} = 27 \nsat^3\, \partial^3_n (E/A)$, and in the isovector
sector the mean-field closed form
%
\begin{equation}
E_{\mathrm{sym}}(\nB) = \frac{k_F^2}{6 E^*_F}
  + \frac{\Gamma_\rho^2(\nB)\, \nB}{2 m_\rho^2},
\qquad L_{\mathrm{sym}} = 3 \nsat\, E_{\mathrm{sym}}'(\nsat),
\qquad K_{\mathrm{sym}} = 9 \nsat^2\, E_{\mathrm{sym}}''(\nsat).
\label{eq:esym}
\end{equation}

The inverse map imposes $\{\nsat, E_{\mathrm{sat}}, m^*/m, K_{\mathrm{sat}},
E_{\mathrm{sym}}, L_{\mathrm{sym}}\}$. The isovector sector closes in two
steps: $\Gamma_\rho(\nsat)$ analytically from Eq.~\eqref{eq:esym}, then
$a_\rho$ from $L_{\mathrm{sym}}$ by a 1-D root. The isoscalar sector has six
couplings $\{\Gamma_\sigma, b_\sigma, c_\sigma, \Gamma_\omega, b_\omega,
c_\omega\}$ against four imposed values, and is closed by two structural
conditions: the cross-constraint $f''_\sigma(1) = f''_\omega(1)$ (which the
published table obeys to $2\times 10^{-3}$), and one shape coefficient pinned
at its published value. Higher derivatives not imposed --- $Q_{\mathrm{sat}}$
and $K_{\mathrm{sym}}$ --- are then \emph{predictions} of the inversion and
are reported as such; imposing $Q_{\mathrm{sat}}$ instead of the pin remains
available as an option (the closure used before this convention).

Two numerical caveats are part of the specification. First,
$Q_{\mathrm{sat}}$ is a third finite difference of a solved quantity; the
$h = 10^{-4}$ stencil carries $\sim 0.1$~MeV of noise, the forward and
inverse maps use the identical stencil so the bias cancels on a round trip,
and the acceptance gate on the isoscalar residual is set by that noise, not
by physics. Second, the isoscalar residual surface has spurious basins;
whether a target inverts can depend on the seed, so the solver retries from
jittered seeds (deterministically) before declaring a target unrepresentable,
and no feasibility boundary may be inferred from a scan run without restarts.

\section{Numerical implementation}
\label{sec:numerics}

Two backends solve the same residual: the reference path (NumPy/SciPy,
finite-difference Jacobian) and the fast path (hand-derived analytic
Jacobian; at $T=0$ residual and Jacobian are Numba-compiled). The two are
held to agreement by backend-parity gates, and the reference path remains the
correctness oracle. At $T>0$ the kinetic derivatives inside the analytic
Jacobian fall back to per-species central differences of the JEL integrals,
which expose no exact derivative. Density sweeps are warm-started --- each
solved point seeds the next, and steps through sharp onsets (hyperon
thresholds, scalar-collapse boundaries) are bisected so the predictor stays
in the corrector's basin. Solved states report residual norms and are gated
at $10^{-10}$; the regression suite pins the symmetric-matter golden point at
$\nB = 0.16\ \mathrm{fm}^{-3}$ and the published DD2 neutron-star properties.

\section{What \texttt{eos\_response} returns}
\label{sec:response}

A second derivative is only defined once one says what is held fixed, so the
conditioning is an explicit argument and each named freeze expands to a
statement about which reactions are faster than the perturbation. Two are
implemented; every other combination raises \texttt{NotImplementedError}
naming itself.

\paragraph{\texttt{frozen='equilibrium'}} Everything re-equilibrates. Wired for
\texttt{beta\_eq\_neutrinoless} only, and computed from the \emph{analytic}
Jacobian rather than by differencing solved points:
%
\begin{center}
\begin{tabular}{lll}
\toprule
key & quantity & availability \\
\midrule
\texttt{cs2\_isothermal} & $c_s^2 = (\partial P/\partial\varepsilon)_T$ along
                   the mode's own sequence & always \\
\texttt{cs2\_adiabatic} & the same at fixed entropy per baryon,
                   $= (C_P/C_V)\,c_{s,\mathrm{iso}}^2$ & always \\
\texttt{chi}     & $\chi_{ab} = \partial n_a/\partial\mu_b$,
                   $a,b \in (B, C, S)$ & always \\
\texttt{C\_V}    & $C_V = (T/\nB)\,(\partial s/\partial T)_{\nB}$ & $T > 0$
                   only \\
\texttt{C\_P}    & $C_P = (T/\nB)\,(\partial s/\partial T)_{P}$ & $T > 0$
                   only \\
\bottomrule
\end{tabular}
\end{center}
%
$C_V$ and $C_P$ are absent at $T = 0$ rather than returned as zero: they are
not defined there, and a zero would be indistinguishable from a computed one.
\texttt{cs2\_adiabatic} is present at $T = 0$ all the same, where $C_P/C_V = 1$
by construction and the two speeds coincide.

\paragraph{\texttt{frozen='composition'}} Every particle fraction held fixed,
i.e.\ reactions slow compared with the perturbation. This is the freeze that
takes a target: the proton fraction \texttt{Y\_p} is a named argument of
\texttt{eos\_response}, and it is a \emph{freeze target}, not one of the
mode's conditions --- which is why it is not among $\nB, T, Y_C, Y_S,
Y_{L_e}$. It returns the same two sound speeds at that proton fraction, by
finite difference along the frozen-$Y_p$ sequence, together with the index
$\Gamma = (\varepsilon + P)/P \; c_{s,\mathrm{iso}}^2$, for nucleonic matter.

\paragraph{What raises.} \texttt{frozen='equilibrium'} in any mode but
\texttt{beta\_eq\_neutrinoless}; the frozen conserved fractions; and the
leptonic re-neutralization variants. Each names itself and points at
\texttt{docs/DEFERRED.md}.

\paragraph{Isothermal against adiabatic.} The conditioning of a second
derivative has two axes here, and each is carried by its own thing, so that no
word does double duty. The \emph{composition} axis is the \texttt{frozen=}
argument: \texttt{frozen='equilibrium'} differentiates along the mode's own
sequence, \texttt{frozen='composition'} at frozen particle fractions. The
\emph{thermal} axis is the key name: \texttt{cs2\_isothermal} holds $T$,
\texttt{cs2\_adiabatic} holds the entropy per baryon, and no quantity in this
package is called a bare \texttt{cs2}. Both freezes return both speeds --- four
numbers, two keys, two arguments.

The word is what makes this necessary. In asteroseismology \cite{ZhaoLattimer2022},
Eq.~(1),
%
\begin{equation}
  \nu_g^2 = g^2 \left( \frac{1}{c_e^2} - \frac{1}{c_s^2} \right)
            e^{\nu - \lambda} ,
\end{equation}
%
``the adiabatic sound speed'' $c_s$ means \emph{frozen composition}, and the
$g$-mode frequency is precisely the difference between it and the equilibrium
$c_e$; given only one of the two, $\nu_g \equiv 0$. In the CompOSE manual
\cite{TypelCompOSE2015} ``adiabatic'' means \emph{fixed entropy}. This library serves
both literatures, so the word is never used unqualified: the composition is
said by the argument and the entropy by the key.

One gap remains, recorded in \texttt{docs/DEFERRED.md}: $\Gamma$ is built on
\texttt{cs2\_isothermal}, so at $T > 0$ it is the isothermal index and the
fixed-entropy one is larger by the same $C_P/C_V$.

\section{Masses}
\label{sec:masses}

The baryon and lepton masses are not DD2 parameters --- they come from the
shared particle table and a parametrization overrides one only through its own
\texttt{Parameters} fields, which for DD2 is the two nucleon masses of
Table~\ref{tab:params}. For completeness, in MeV:
%
\begin{center}
\begin{tabular}{lrl@{\qquad}lrl}
\toprule
species & mass & $g_i$ & species & mass & $g_i$ \\
\midrule
$n$        & $939.56536$  & $2$ & $\Xi^-$      & $1321.71$   & $2$ \\
$p$        & $938.27203$  & $2$ & $\Xi^0$      & $1314.86$   & $2$ \\
$\Lambda$  & $1115.683$   & $2$ & $\Delta$ quartet & $1232.00$ & $4$ \\
$\Sigma^-$ & $1197.449$   & $2$ & $e^-$        & $0.510999$  & $2$ \\
$\Sigma^0$ & $1192.642$   & $2$ & $\mu^-$      & $105.6584$  & $2$ \\
$\Sigma^+$ & $1189.370$   & $2$ & & & \\
\bottomrule
\end{tabular}
\end{center}
%
The $\Delta$ degeneracy is $4$ because the quartet member carries spin $3/2$;
every octet baryon and both leptons carry $2$. The meson-gas masses and
degeneracies are in Sec.~\ref{sec:gas}.

\bibliography{../../docs/eos}

\end{document}
